% ============================================================
\chapter{Current Research Directions}
\label{ch:research}
% ============================================================

\begin{intuition}
Geometric deep learning is a rapidly expanding field. This chapter
presents the most active research directions and open problems that
will shape the next advances in the field.
\end{intuition}

% ============================================================
\section{Geometric foundation models}
\label{sec:foundation_models}
% ============================================================

\begin{definition}[Geometric foundation model]
\index{foundation model}
A \textbf{geometric foundation model} is a large-scale pretrained
model on structured data (graphs, molecules, point clouds) that is
transferable to multiple downstream tasks via fine-tuning or
prompting.
\end{definition}

\begin{example}[Examples of geometric foundation models]
\begin{itemize}
  \item \textbf{GEM} (Fang et al.\ 2022): foundation model for
        molecules, pretrained on 20M structures.
  \item \textbf{Uni-Mol} (Zhou et al.\ 2023): unified 2D/3D
        molecular representation.
  \item \textbf{Point-MAE} (Pang et al.\ 2022): masked autoencoder
        for 3D point clouds.
\end{itemize}
\end{example}

\begin{remark}
Unlike text foundation models (GPT, Claude), geometric foundation
models must respect data \textbf{symmetries} (permutation invariance,
rotation equivariance). This imposes strong architectural constraints.
\end{remark}

\begin{proposition}[Scaling challenges]
Scaling GNNs poses specific problems:
\begin{enumerate}
  \item \textbf{Neighborhood explosion}: $k$ layers imply
        $\mathcal{O}(d^k)$ neighbors ($d$ = average degree).
  \item \textbf{Over-smoothing}: representations converge with
        depth.
  \item \textbf{Over-squashing}: long-range information is
        compressed through bottlenecks.
  \item \textbf{Variable-size graphs}: batching is non-trivial.
\end{enumerate}
\end{proposition}

% ============================================================
\section{Equivariant transformers}
\label{sec:equivariant_transformers}
% ============================================================

\begin{definition}[Geometric transformer]
\index{geometric transformer}
A \textbf{geometric transformer} is a transformer whose attention
mechanism integrates geometric information while respecting
symmetries:
\[
  \alpha_{ij} = \frac{\exp\!\left(\frac{q_i^\top k_j}{\sqrt{d}}
  + b(x_i, x_j)\right)}{\sum_{l} \exp\!\left(
  \frac{q_i^\top k_l}{\sqrt{d}} + b(x_i, x_l)\right)}
\]
where $b(x_i, x_j)$ is a geometry-dependent attention bias
(distance, angle, relative orientation).
\end{definition}

\begin{example}[Example architectures]
\begin{itemize}
  \item \textbf{SE(3)-Transformer} (Fuchs et al.\ 2020):
        $\mathrm{SE}(3)$-equivariant attention with tensor features.
  \item \textbf{Equiformer} (Liao \& Smidt, 2023): equivariant
        transformer with tensor products in attention.
  \item \textbf{GPS} (Rampa\v{s}ek et al.\ 2022): combines local
        message passing and global attention.
\end{itemize}
\end{example}

\begin{theorem}[Expressiveness of Graph Transformers]
\index{Graph Transformer}
A Graph Transformer with spectral positional encoding and global
attention is \textbf{strictly more expressive} than the $1$-WL test.
With $\mathrm{SE}(3)$-equivariant positional encodings, it can
approximate any continuous equivariant function.
\end{theorem}

\begin{keyformulas}[title=\textbf{Equivariant attention}]
\[
  h_i' = \sum_{j} \alpha_{ij} \cdot V(h_j, x_j - x_i)
\]
where $V : \R^d \times \R^3 \to \R^{d'}$ is a value function
depending on relative position, and $\alpha_{ij}$ is
$\mathrm{SE}(3)$-invariant.
\end{keyformulas}

% ============================================================
\section{Diffusion models on manifolds}
\label{sec:diffusion_manifolds}
% ============================================================

\begin{definition}[Riemannian diffusion]
\index{Riemannian diffusion}
A \textbf{Riemannian diffusion model} (De Bortoli et al.\ 2022)
generalizes diffusion models (DDPM) to Riemannian manifolds
$(\mathcal{M}, g)$:
\begin{enumerate}
  \item \textbf{Forward process}: Riemannian Brownian motion
        $dX_t = \sqrt{2\beta_t}\, dB_t^{\mathcal{M}}$.
  \item \textbf{Score matching}: learn the score
        $\nabla_{\mathcal{M}} \log p_t(x)$ on the manifold.
  \item \textbf{Reverse process}: integrate the reverse equation
        on $\mathcal{M}$ with the learned score.
\end{enumerate}
\end{definition}

\begin{example}[Applications of geometric diffusion]
\begin{itemize}
  \item \textbf{3D molecule generation}: diffusion on $\R^{3n}$
        with $\mathrm{SE}(3)$ equivariance
        (EDM, Hoogeboom et al.\ 2022).
  \item \textbf{Conformation generation}: diffusion on torsion
        angles (toric manifold).
  \item \textbf{Protein generation}: diffusion on
        $\mathrm{SO}(3)^n$ for residue frames
        (RFDiffusion, Watson et al.\ 2023).
\end{itemize}
\end{example}

\begin{attention}[title=\textbf{Technical challenges}]
Diffusion on manifolds requires:
\begin{itemize}
  \item The exponential map $\exp_x : T_x\mathcal{M} \to
        \mathcal{M}$ for integration steps.
  \item The Riemannian logarithm $\log_x : \mathcal{M} \to
        T_x\mathcal{M}$ for the score.
  \item Parallel transport to move vectors between tangent spaces.
\end{itemize}
\end{attention}

% ============================================================
\section{Geometric self-supervised learning}
\label{sec:geometric_ssl}
% ============================================================

\begin{definition}[Self-supervised learning on graphs]
\index{self-supervised learning}
\textbf{Self-supervised learning} (SSL) on graphs learns
representations without labels via auxiliary tasks:
\begin{enumerate}
  \item \textbf{Contrastive}: maximize similarity between two
        augmented views of the same graph (GraphCL, GCA).
  \item \textbf{Predictive}: predict masked properties (node
        features, edges, subgraphs).
  \item \textbf{Generative}: reconstruct the graph from a
        compressed representation (GraphMAE).
\end{enumerate}
\end{definition}

\begin{proposition}[Geometric augmentations]
Augmentations for graphs include:
\begin{itemize}
  \item Random node or edge dropout.
  \item Feature perturbation.
  \item Random subgraph extraction.
  \item Heat diffusion (feature smoothing).
\end{itemize}
For 3D data, add: rotation, translation, Gaussian noise, partial
sampling.
\end{proposition}

\begin{theorem}[Contrastive SSL guarantees]
\index{contrastive learning}
Under regularity assumptions on the data distribution and
augmentations, contrastive representations are \textbf{linearly
separable} for downstream tasks, with an error bound controlled
by the alignment and uniformity of the representations.
\end{theorem}

% ============================================================
\section{Open problems}
\label{sec:open_problems}
% ============================================================

\begin{remark}[Major open problems]
\index{open problems}
\begin{enumerate}
  \item \textbf{Optimal expressiveness}: what is the minimal GNN
        architecture to match $k$-WL for a given $k$?
  \item \textbf{Over-squashing}: how to enable long-range
        information flow without expensive global attention
        ($\mathcal{O}(n^2)$)?
  \item \textbf{Out-of-distribution generalization}: do GNNs
        generalize to graphs with very different structures from
        training?
  \item \textbf{Approximate equivariance}: how to handle
        approximate symmetries (quasi-symmetries) of real-world
        data?
  \item \textbf{Interpretability}: how to explain GNN decisions
        in terms of meaningful substructures?
  \item \textbf{Scalability}: how to train GNNs on graphs with
        billions of nodes?
\end{enumerate}
\end{remark}

\begin{definition}[Over-squashing]
\index{over-squashing}
\textbf{Over-squashing} refers to the phenomenon where information
from distant nodes is exponentially compressed when passing through
message-passing layers. Formally, the Jacobian:
\[
  \left\|\frac{\partial h_v^{(L)}}{\partial h_u^{(0)}}\right\|
  \leq C \cdot \lambda_{\max}^L \cdot
  \prod_{\ell=1}^L \mathrm{Lip}(\sigma_\ell)
\]
decays exponentially with graph distance if $\lambda_{\max} < 1$.
\end{definition}

% ============================================================
\section{Exercises}
% ============================================================

\begin{exercise}
Explain why Laplacian positional encodings are not unique (sign,
ordering). Propose a solution to make them invariant.
\end{exercise}

\begin{exercise}
For a molecular graph, propose three contrastive augmentations
that preserve chemical identity and three that do not. Justify.
\end{exercise}

\begin{exercise}[Over-squashing]
Compute the bound on $\left\|\frac{\partial h_v^{(L)}}
{\partial h_u^{(0)}}\right\|$ for a path graph of length $L$
with a GCN with unit weights and ReLU activation.
\end{exercise}

\begin{exercise}
Show that Brownian motion on the sphere $S^2$ converges to the
uniform distribution. What is the convergence rate as a function
of the first nonzero eigenvalue of the Laplacian?
\end{exercise}

\begin{exercise}[Project]
Implement a simple equivariant diffusion model to generate 3D point
clouds (e.g., shapes from ShapeNet). Evaluate quality using Chamfer
distance and Betti numbers.
\end{exercise}
